% =====================================================================
\lezione{12/10/2026}{Debrief del Lab~1; introduzione a ISO/IEC~27001; presentazione del caso Sylvane}
  {Prof.~Giuseppe Vaciago; Gianluca Gilardi}
  {presentazioni dei gruppi sul caso Arvenza e commenti; triade CIA; esiti della gestione del rischio
   (mitigare, accettare, trasferire); ISMS e sue componenti; famiglia ISO 27k; presentazione del
   Lab~2 (Sylvane Retraite)}
\label{lez:2026-10-12}
% =====================================================================

\begin{nota}[Fonte]
Trascrizione automatica della registrazione, revisionata (pausa tra il minuto 66 e l'87 circa). La
parte delle presentazioni è in larga parte inintelligibile: si riportano solo i commenti dei docenti
e i contenuti ricorrenti. Slide di riferimento: \emph{04.~polito\_iso\_2026} (Fundamentals and
ISO/IEC Standards). Il caso del Lab~2 è nel Capitolo~\ref{lab:02}.
\end{nota}

\section{Sintesi}
La lezione ha tre parti. Nei primi 50 minuti alcuni gruppi presentano il proprio lavoro sul caso
Arvenza (Lab~1) e i docenti commentano. Segue la prima lezione di \textbf{Gianluca Gilardi}:
introduzione all'\textbf{ISO/IEC~27001} come sistema di gestione basato sul rischio, con triade CIA,
tre esiti della gestione del rischio e componenti dell'ISMS. Infine Gianluca presenta il caso
\textbf{Sylvane Retraite} per il laboratorio di venerdì~16/10.

% ---------------------------------------------------------------------
\section{Organizzazione}
\begin{itemize}
  \item \textbf{Esiti del Lab~1}: ogni gruppo ha ricevuto una valutazione via email. Soglie
    comunicate: sotto 20 nessun punto bonus; da 20 a 25 si ottiene il punto bonus; da 25 a 30 si può
    anche presentare il caso in aula. Chi ritiene di avere diritto al bonus ma non è in elenco deve
    segnalarlo al docente.
  \item \textbf{Ritardi}: tollerati i lavori arrivati entro le 14:00 e quelli inviati in tempo ma
    senza allegato (``siete esperti di sicurezza: ricordatevi l'allegato''); non accettati quelli
    arrivati il giorno dopo.
  \item \textbf{Regole formali per i prossimi lab}: nomi di tutti i componenti in prima pagina e
    \textbf{file nominato con il nome del gruppo}. Con oltre 90 studenti e più di 25 gruppi, un lavoro
    non attribuibile viene scartato.
  \item \textbf{Consegna del Lab~2}: caricamento sul repository SharePoint indicato nel caso
    (preferito); in alternativa email al docente o sezione \emph{Elaborati} del portale, ma è meglio un
    solo punto di raccolta.
  \item Su Zoom non si useranno più le sale separate: le domande si fanno nella sala principale (o in
    chat se brevi).
  \item Le slide ISO vengono caricate sul portale la sera stessa. Agli studenti senza gruppo viene
    chiesto di trovarne uno (anche online).
  \item Il materiale del caso Sylvane è online dalla mattina; si può studiare in anticipo, ma il lavoro
    si svolge venerdì.
\end{itemize}

% ---------------------------------------------------------------------
\section{Debrief del Lab~1 (Arvenza)}

\subsection{Che cosa è emerso dalle presentazioni}
Le presentazioni, da cinque--sette minuti, hanno mostrato approcci diversi, tutti basati sulla
formula ENISA:
\begin{itemize}
  \item \textbf{EI = 1} per tutti i gruppi di dati (nome, codice fiscale, IBAN identificano
    direttamente la persona);
  \item \textbf{intento malevolo +0,5} ovunque (l'attaccante voleva deviare gli stipendi);
  \item due \textbf{perimetri}: solo ciò che risulta letto dai log, oppure l'intera mailbox;
  \item più gruppi hanno sottolineato che \textbf{la severità non è il rischio}: alcuni dati hanno
    severità alta ma nessuna prova di accesso (rischio basso), altri severità media ma rischio molto
    alto (IBAN già usati per la frode);
  \item date di inizio delle 72 ore diverse: 20/9 (alert di sicurezza) oppure 28/9 (segnalazione del
    payroll bureau);
  \item decisioni variabili: quasi tutti hanno respinto la posizione del CEO e combinato le altre due.
\end{itemize}

\subsection{Commenti di Gianluca}
\begin{esame}
\begin{itemize}
  \item \textbf{Certificati di malattia}: molti li hanno valutati DPC~2 senza motivare. Gianluca è
    personalmente d'accordo, ma ai sensi dell'\textbf{art.~9 GDPR} ogni dato relativo alla salute ha
    la stessa natura e ricade nel \textbf{DPC~4} del modello ENISA. Abbassarlo è possibile solo
    \textbf{argomentando}.
  \item \textbf{Data di consapevolezza}: un ispettore del Garante contesterebbe quasi certamente una
    data di consapevolezza così tardiva. La violazione era forse attribuibile già all'inizio della
    cronologia (richieste di modifica dell'IBAN respinte e verificate con gli interessati). Chi
    sostiene una data tardiva deve essere pronto a dimostrare perché la conoscenza effettiva è
    arrivata dopo.
  \item In generale: si può sostenere quasi tutto, ma ogni affermazione e decisione deve avere
    \textbf{fondamento ed evidenze}. Nel caso successivo servirà un'analisi ancora più accurata.
\end{itemize}
\end{esame}

\paragraph{Commenti del docente.} Si è apprezzata la divisione dei ruoli nel gruppo; in una
presentazione devono parlare tutti i componenti, non uno solo. Quando si parla in pubblico si guarda
chi \emph{non} sta ascoltando, non il docente. Un gruppo ha riconosciuto di aver sovrastimato il CB
aggiungendo integrità e disponibilità ``per prudenza''.

% ---------------------------------------------------------------------
\section{ISO/IEC 27001: introduzione (Gianluca Gilardi)}

\begin{nota}[Il docente]
Gianluca è stato avvocato per 13 anni, poi \emph{digital forensics} per otto anni (spiegava la
tecnologia ad avvocati e giudici e i requisiti legali ai tecnici); da otto anni è CEO della società di
legal tech fondata dal prof.~Vaciago. Assiste aziende durante data breach e audit.
\end{nota}

\subsection{Perché ISO 27001}
\begin{itemize}
  \item La certificazione ISO/IEC 27001 significa che una \textbf{terza parte fidata} ha verificato
    che l'organizzazione ha misure e controlli appropriati per la sicurezza delle informazioni.
  \item \textbf{Non è richiesta da nessuna legge}: è uno schema volontario. Dal punto di vista del
    mercato, però, conoscerla è quasi un requisito per chi lavora in cybersecurity, anche senza
    certificarsi.
  \item È un sistema \textbf{basato sul rischio}: i controlli si adottano dove il rischio valutato lo
    richiede, per portarlo a un livello accettabile.
  \item Esistono alternative: negli USA i clienti chiedono più spesso il \textbf{SOC~2}; in Europa è
    più accettata la ISO. La differenza è di mercato, non tecnica.
  \item ``Sono troppo piccolo'': nel 2026 nessuna organizzazione è esente, perché tutte raccolgono,
    trattano, conservano e trasmettono informazioni che sono \textbf{asset} esposti a rischio.
    Battuta del settore: non chiederti \emph{se} subirai un incidente, ma \emph{quando}.
\end{itemize}

\subsection{I tre esiti della gestione del rischio}
Secondo Gianluca ogni processo di gestione del rischio si chiude con una di tre decisioni:
\begin{enumerate}
  \item \textbf{mitigare}: adottare controlli che riducano il rischio a un livello accettabile;
  \item \textbf{accettare} il rischio (mitigato o no);
  \item \textbf{trasferire} il rischio a un terzo (il ``Santo Graal'' del risk manager), come con
    l'assicurazione auto. Bisogna però poter dimostrare che la scelta del terzo è solida: non basta
    una firma su un foglio.
\end{enumerate}

\begin{nota}[Dalle slide successive del corso]
Le slide \emph{06.~27001\_2026} elencano anche una quarta opzione, \textbf{evitare} (cessare
l'attività che genera il rischio), che corrisponde allo ``stop del trattamento'' visto per la DPIA
(Capitolo~\ref{lez:2026-10-05}).
\end{nota}

\subsection{La triade CIA}
Dalle slide (definizioni ISO 27000): la sicurezza delle informazioni è la preservazione di
\textbf{riservatezza, integrità e disponibilità}; possono rilevare anche autenticità,
accountability, non ripudio e affidabilità.
\begin{itemize}
  \item \textbf{Riservatezza}: l'informazione non è resa disponibile né divulgata a persone, entità o
    processi non autorizzati. Non tutte le informazioni hanno la stessa sensibilità:
    \textbf{classificare} serve a non sprecare risorse su informazioni irrilevanti.
  \item \textbf{Integrità}: accuratezza e completezza.
  \item \textbf{Disponibilità}: accessibile e utilizzabile su richiesta da un'entità autorizzata.
\end{itemize}

\begin{esame}
\begin{itemize}
  \item Le \textbf{definizioni formali} di riservatezza, integrità e disponibilità sono molto
    importanti ed è probabile che vengano chieste all'esame.
  \item Anche una \textbf{perdita temporanea di disponibilità} di dati personali (poche ore) è un
    data breach. Il severity assessment potrà concludere che il rischio è così basso da non
    richiedere notifiche, ma la violazione va comunque \textbf{annotata nel registro}.
\end{itemize}
\end{esame}

\subsection{Che cos'è un ISMS}
Un \textbf{Information Security Management System} è un approccio sistematico per stabilire,
attuare, gestire, monitorare, riesaminare, mantenere e migliorare la sicurezza delle informazioni
di un'organizzazione. Un \emph{sistema}, non un patchwork: componenti tecniche, organizzative e
documentali devono stare in armonia (approccio olistico).

\paragraph{Monitoraggio e data governance.} Nel Lab~1 era difficile stabilire quando il titolare
fosse venuto a conoscenza della violazione: è il problema di chi non monitora. Il peccato più grave
nella gestione dei dati è non saper provare \emph{cosa}, \emph{dove}, \emph{quando}, \emph{chi} e
\emph{perché}: prima ancora della sicurezza viene la \textbf{data governance}.

\paragraph{Componenti dell'ISMS (slide).}
\begin{enumerate}[label=\alph*)]
  \item \textbf{policy}: la parte più facile da scrivere (spesso quasi un copia-incolla dei
    controlli), più difficile da mantenere aggiornata;
  \item \textbf{persone con responsabilità definite}: senza responsabilità formali (chi monitora, chi
    segnala, chi decide) non c'è un sistema di gestione ma un ``sistema della speranza'';
  \item \textbf{processi di gestione}: definizione delle policy, consapevolezza e competenze
    (gli auditor intervistano anche chi sta alla reception), pianificazione, attuazione, operatività,
    valutazione delle prestazioni, \textbf{riesame della direzione}, miglioramento;
  \item \textbf{informazioni documentate}: niente ``tradizione orale'';
  \item \textbf{valutazione del rischio} per la sicurezza delle informazioni;
  \item \textbf{trattamento del rischio}, con scelta e attuazione dei controlli.
\end{enumerate}
Le componenti a)--d) sono comuni a qualsiasi sistema di gestione; e) e f) sono specifiche
dell'ISMS.

\begin{itemize}
  \item \textbf{Ruolo della direzione}: le decisioni strategiche sulla sicurezza vengono dall'alto.
    Uno dei primi controlli in audit è come il management è coinvolto e informato.
  \item \textbf{Miglioramento continuo}: un ISMS nasce e cresce. Se l'auditor non vede miglioramenti
    solleva un rilievo, che può essere un'\emph{opportunità di miglioramento}, una \emph{non
    conformità minore} o una \emph{non conformità maggiore}.
\end{itemize}

\subsection{Valutare e trattare il rischio}
\begin{itemize}
  \item Rischio = severità (o impatto) $\times$ probabilità. Dalle slide: rischio come ``effetto
    dell'incertezza sugli obiettivi''; il rischio informativo è la possibilità che una minaccia sfrutti
    una vulnerabilità di un asset.
  \item Il risk manager prudente valuta sempre lo \textbf{scenario peggiore}, compresi i
    \textbf{cigni neri}: il COVID-19 è stato il più duro esercizio di gestione del rischio, e il
    mondo non era pronto perché la probabilità era stata giudicata troppo bassa.
  \item Il management chiederà di tagliare i costi: bisogna \textbf{prioritizzare} in base a
    criticità degli asset, impatto, probabilità e \textbf{rapporto costi/benefici} dei controlli.
  \item Esistono \emph{quick win}: i \textbf{backup}. Gianluca cita un cliente bloccato oltre due mesi
    da un ransomware perché non aveva backup dei dati cifrati.
  \item Alcuni controlli sono specifici per rischi ristretti, altri hanno portata più ampia: anche
    questo incide sul costo.
\end{itemize}

\subsection{Benefici e famiglia ISO 27k}
\begin{itemize}
  \item Ridurre il rischio per la sicurezza delle informazioni significa anche \textbf{ridurre il
    rischio di compliance}: è un investimento in conformità (le sanzioni sono rilevanti).
  \item La ISO 27001 è lo standard più usato per attuare e certificare un ISMS. La famiglia 27k conta
    \textbf{oltre 70 standard} (cloud, protezione dei dati personali, ecc.).
  \item Secondo Gianluca le altre certificazioni 27k sono ``figlie'' della 27001: senza la 27001 non
    si possono ottenere.
\end{itemize}

\begin{integrazione}[27k ``figlie'' della 27001]
Era vero per la ISO/IEC~27701:2019 (estensione della 27001). Dall'edizione 2025 la 27701 è uno standard
autonomo, certificabile anche senza 27001 (vedi Capitolo~\ref{cap:iso}).
\end{integrazione}

\begin{esame}
Il docente sottolinea che l'\textbf{ISMS è la base della ISO 27001}: è il primo concetto da fissare.
\end{esame}

% ---------------------------------------------------------------------
\section{Presentazione del Lab~2: Sylvane Retraite (Gianluca)}
Nuovo caso di data breach, più complesso sotto certi aspetti e più semplice sotto altri (dati meno
eterogenei, cronologia più breve). Dettagli e svolgimento nel Capitolo~\ref{lab:02}.
\begin{itemize}
  \item \textbf{Wall clock} 105 minuti (tempo reale del lab) e \textbf{exercise clock} tre giorni (il
    tempo della storia, nei tre giorni dopo la presa di conoscenza).
  \item La violazione arriva da una \textbf{comunicazione del responsabile del trattamento}: questo
    semplifica l'individuazione del momento in cui parte l'orologio per il titolare.
  \item Il responsabile ha già notificato l'incidente a un'autorità di settore \textbf{ai sensi del
    DORA}: diverse norme, diverse notifiche.
  \item Stavolta il \textbf{modulo di severity assessment è obbligatorio}; se si decide di notificare,
    si compila anche il file Excel del modulo di notifica (modello usato da molte autorità europee), e
    per gli interessati va scritto il testo effettivo della comunicazione (art.~34).
  \item Durante il lab arriveranno \textbf{curveball}, nuove informazioni che possono cambiare la
    valutazione: ``la vita è ciò che accade mentre fai l'esercizio''.
  \item \textbf{Honne e tatemae}: il modulo chiede sia la giustificazione da dare all'autorità
    (\emph{tatemae}, la facciata) sia il vero ragionamento interno (\emph{honne}, la verità, anche
    ``temo di essere licenziato'').
\end{itemize}

% ---------------------------------------------------------------------
\section{Domande aperte e punti da verificare}
\begin{verifica}
\begin{itemize}
  \item \textbf{Cognome di Gianluca}: ``Gilardi'' (dalla trascrizione del 09/10 e dall'indirizzo del
    repository SharePoint del Lab~2). La trascrizione oscilla tra ``Gilardi'' e ``Gillardi''.
  \item Soglie del bonus (20 / 25 / 30): dal parlato del docente; confermare con l'email ricevuta.
  \item ``Senza 27001 non si ottengono le altre certificazioni 27k'': affermazione di Gianluca, vera
    per le estensioni come la 27701; da verificare per gli altri standard della famiglia.
  \item Gianluca indica come \emph{honne} il ``paragrafo 10'': nel modulo il paragrafo~9 è la
    giustificazione (tatemae) e il 10 la strategia interna (honne). Coerente, ma verificare.
\end{itemize}
\end{verifica}
